\documentclass[11pt]{ctexart}
\usepackage[a4paper,margin=1in]{geometry}
\usepackage{graphicx}
\usepackage{xcolor}
\usepackage{hyperref}
\definecolor{refgray}{gray}{0.45}
\title{\textbf{Market Views｜全球投行叙事汇编}\\[0.4em]{\large AI系统级竞争加速，亚洲信贷扩张至18年高位，消费与零售结构分化加剧}}
\author{KC桌面 自动生成}
\date{260724}
\begin{document}
\maketitle
\tableofcontents
\newpage
\section{一页摘要}
\begin{itemize}
  \item 亚洲（除中国）信贷增速创18年新高，由AI、能源和供应链资本支出驱动，信贷扩张广度持续扩大。
  \item AI产业从单芯片算力竞赛转向系统级互联与集群效率竞争，中国厂商在WAIC 2026展示千卡级SuperPod方案，但软件层AI货币化仍处早期。
  \item 欧元区通胀回落受能源价格反弹干扰，但第二轮效应（工资-通胀螺旋）尚未显现；日本央行7月会议料维持利率不变。
  \item 中国快递物流业从价格战转向质量驱动，单价回升至7.77元，头部企业盈利拐点确认；美国运动服饰龙头Lululemon美洲业务规模或高估25\%-60\%。
  \item 模拟芯片行业周期性复苏加速，工业与汽车终端市场同步改善；全球电网投资结构性增长，数据中心订单陡峭增长。
  \item 中国医疗健康板块受AI资金外溢与CXO结构性增持推动反弹，但政策不确定性与资金持续性存分歧。
  \item 中国连锁药房同店增长因中小药店出清而加速，驱动力为客流增长；全球光纤价格因AI数据中心需求创历史新高，但中国产能扩张预示2026年或为周期顶点。
  \item 美国散户资金从被动配置转向主动交易，动量因子因极端拥挤回撤约14\%，科技ETF流入转负但个股净买入持续。
\end{itemize}
\section{全球投行叙事汇编}
今日纳入 40 篇投行/券商研究，按 7 个市场、资产与产业主题系统整理。
\newpage
\subsection{宏观与利率：全球资金流向与通胀传导的再定价}
\textbf{全球宏观环境正经历多重力量交织：亚洲（除中国）信贷扩张至18年高位，由AI、能源和供应链资本支出驱动；欧元区通胀前景受能源价格反弹与需求走弱的双向拉扯，但第二轮效应尚未显现；日本央行维持观望，但GPIF潜在的本土资产配置调整可能对日元和全球资金流产生深远影响。}
\subsubsection*{主流共识}
\begin{itemize}
  \item 亚洲（除中国）信贷增速创18年新高，主要由企业资本支出和贸易活动驱动，且信贷扩张广度正在扩大。
  \item 欧元区通胀回落进程受到能源价格反弹干扰，但第二轮效应（工资-通胀螺旋）尚未出现。
  \item 日本央行7月会议将维持利率不变，下一次调整更可能在2027年1月。
\end{itemize}
\subsubsection*{分歧与条件差异}
\begin{itemize}
  \item 关于AI货币化能否抵消传统业务拖累：MS认为Salesforce的AI收入增速虽快，但体量尚不足以覆盖Commerce和Tableau的结构性下滑，有机增长拐点未至。
  \item 关于中国房地产市场复苏形态：HSBC认为二手房流动性改善和高端需求韧性支撑K型复苏，但市场对盈利修复的时间点（2027年）存在分歧。
  \item 关于日本GPIF调整资产配置的节奏与影响：JPM测算33.8万亿日元回流可能推动USD/JPY下跌约15日元，但实际影响取决于资金流入是渐进式还是集中释放。
\end{itemize}
\subsubsection*{机构观点}
\paragraph{MS} 亚洲（除中国）银行信贷增速升至8.5\%（2026年5月），为18年最高，由AI基础设施、能源、国防和供应链资本支出超级周期驱动。资本品进口增速达33\%（3MMA），PPI升至9.0\%，企业贷款增速同步创18年新高。家庭贷款随就业改善从5.6\%回升至6.6\%，信贷扩张广度扩大。中国信贷增速持续放缓，形成鲜明对比。
{\small\color{refgray}数据：亚洲（除中国）信贷增速8.5\%，18年最高\par}
{\small\color{refgray}数据：资本品进口增速33\%（3MMA），2004年以来最高\par}
{\small\color{refgray}数据：PPI升至9.0\%，四年高点\par}
{\small\color{refgray}数据：家庭贷款增速从5.6\%回升至6.6\%\par}
{\small\color{refgray}边际变化：信贷扩张从科技领域扩散至非科技领域，PPI上升对营运资金需求的拉动成为新中间变量。\par}
\paragraph{MS} Salesforce的Agentforce年化经常性收入突破10亿美元，增速超200\%，token消耗量环比增长152\%。但AI势头尚未转化为有机增长拐点，cRPO连续两个季度未超预期。Commerce和Tableau板块（合计约30\%收入）受Shopify竞争和AI替代不确定性拖累，传统业务的结构性影响仍在持续。
{\small\color{refgray}数据：Agentforce ARR突破10亿美元，增速超200\%\par}
{\small\color{refgray}数据：Token消耗量环比增长152\%\par}
{\small\color{refgray}数据：Commerce和Tableau合计约占FY26收入的30\%\par}
{\small\color{refgray}数据：cRPO连续两个季度未超预期\par}
{\small\color{refgray}边际变化：AI货币化仍处早期，传统业务拖累抵消AI增长动能，有机增长拐点未至。\par}
\paragraph{NOM} USD/CNY中间价模型投影值为6.7716，较实际中间价6.7933低217个基点，计入逆周期因子后投影值为6.7770，仍低163个基点。模型误差扩大表明央行在中间价定价中引入额外稳定因子，而非完全跟随市场波动。逆周期因子贡献约54个基点调整。
{\small\color{refgray}数据：模型投影中间价6.7716，实际6.7933\par}
{\small\color{refgray}数据：模型误差217个基点\par}
{\small\color{refgray}数据：逆周期因子调整约54个基点\par}
{\small\color{refgray}数据：计入逆周期因子后投影6.7770\par}
{\small\color{refgray}边际变化：模型误差扩大，央行中间价管理意图量化信号增强。\par}
\paragraph{GS} 欧元区通胀前景受两股力量拉扯：物价数据和调查走弱，但中东冲突与俄罗斯炼厂停产推高能源价格。BVAR模型显示能源冲击的直接效应强劲，间接效应在4月累积但6月已消退，未出现第二轮效应。能源期货价格回升将通胀预测拉回欧洲央行基线。
{\small\color{refgray}数据：BVAR模型包含19个变量\par}
{\small\color{refgray}数据：直接效应在3月强劲，间接效应6月消退\par}
{\small\color{refgray}数据：未出现第二轮效应（工资-通胀螺旋）\par}
{\small\color{refgray}数据：能源期货价格回升使通胀预测与ECB基线一致\par}
{\small\color{refgray}边际变化：能源价格上涨抵消了此前物价调查和通胀数据所显示的进展，但第二轮效应尚未出现。\par}
\paragraph{GS} 日本央行7月会议将维持政策不变，下一次利率调整更可能在2027年1月。6月短观调查显示企业信心良好，投资意愿积极。原油价格低于4月展望报告，央行可能上调增长预测并小幅下调通胀预测。但调整时机的概率分布已偏向前移，不排除10月调整。
{\small\color{refgray}数据：7月会议预计维持利率不变\par}
{\small\color{refgray}数据：下一次调整更可能在2027年1月\par}
{\small\color{refgray}数据：布伦特原油价格低于4月展望报告水平\par}
{\small\color{refgray}数据：春季劳资谈判实现约3\%基础工资涨幅\par}
{\small\color{refgray}边际变化：调整时机概率分布偏向前移，10月调整可能性上升。\par}
\paragraph{HSBC} 中国房地产市场呈现K型复苏：10城二手房成交量同比上升约7\%，上海、北京、苏州6月同比增幅分别达21\%、11\%、28\%。高端需求保持韧性，杭州上半年录得创纪录豪宅交易量。越秀地产盈利预警（核心净利润同比降90-95\%）未引发市场波动，表明业绩压力已被消化，市场关注转向2027年盈利修复。
{\small\color{refgray}数据：10城二手房成交量同比+7\%\par}
{\small\color{refgray}数据：上海6月二手房成交量同比+21\%\par}
{\small\color{refgray}数据：越秀地产核心净利润同比降90-95\%\par}
{\small\color{refgray}数据：杭州上半年豪宅交易量创纪录\par}
{\small\color{refgray}边际变化：二手房流动性改善和高端需求韧性支撑K型复苏，市场预期已提前反映2027年盈利修复。\par}
\paragraph{JPM} 日本GPIF若仅将国内债券和股票配置调整至现有目标区间上限（各31\%），可产生约33.8万亿日元买盘/外币卖盘，理论上推动USD/JPY下跌约15日元。但实际影响取决于资金流入节奏，更可能是渐进式。其他日本机构投资者（如生命保险公司持有约100万亿日元海外资产）的跟进意愿是关键变量。
{\small\color{refgray}数据：GPIF国内债券配置26.91\%，目标上限31\%\par}
{\small\color{refgray}数据：GPIF国内股票配置23.81\%，目标上限31\%\par}
{\small\color{refgray}数据：潜在日元买盘约33.8万亿日元\par}
{\small\color{refgray}数据：理论USD/JPY影响约15日元\par}
{\small\color{refgray}边际变化：GPIF调整资产配置的讨论升温，潜在资金回流规模巨大，但节奏和机构跟进存在不确定性。\par}
\subsubsection*{关键数据}
\begin{itemize}
  \item 亚洲（除中国）信贷增速8.5\%，18年最高
  \item 资本品进口增速33\%（3MMA），2004年以来最高
  \item PPI升至9.0\%，四年高点
  \item Agentforce ARR突破10亿美元，增速超200\%
  \item USD/CNY中间价模型误差217个基点
  \item GPIF潜在日元买盘约33.8万亿日元
  \item 10城二手房成交量同比+7\%
  \item 欧元区能源冲击未出现第二轮效应
\end{itemize}
\begin{figure}[htbp]
\centering
\includegraphics[width=0.88\linewidth]{figures/fig_018.jpg}
\caption{Exhibit 1 - Exhibit 1: Credit growth in Asia ex China at a 18-year high Bank credit is still the most important}
\end{figure}
\begin{figure}[htbp]
\centering
\includegraphics[width=0.88\linewidth]{figures/fig_030.jpg}
\caption{Exhibit 1 - Exhibit 2: We Use a Bayesian Vector Auto-Regression to Jointly Model 19 Variables}
\end{figure}
\begin{figure}[htbp]
\centering
\includegraphics[width=0.88\linewidth]{figures/fig_047.jpg}
\caption{Exhibit 1 - Exhibit 1: Business Sentiment Is Favorable BOJ Tankan, Business Conditions DI (All Enterprise Sizes) \#\# Exhibit 2: Companies Maintain a Positive Investment Stance BOJ Tankan, Capex Plans (All Enterprises, Excludes Land, Includes}
\end{figure}
\begin{figure}[htbp]
\centering
\includegraphics[width=0.88\linewidth]{figures/fig_078.jpg}
\caption{Figure 2 - Figure 2: Cumulative flow into foreign bonds trillion yen No data for Pension Fund pre 2014 Figure 3: Cumulative flow into foreign equities trillion yen No data for Pension Fund pre 2014 Figure 4: \$30\textbackslash{}\%\$ of lifers' security as}
\end{figure}
{\small\color{refgray}本节综合 7 篇研究；涉及 MS、NOM、GS、HSBC、JPM。}
\newpage
\subsection{AI基础设施与软件变现：系统级竞争加速，货币化路径分化}
\textbf{AI产业正从单芯片算力竞赛转向系统级互联与集群效率竞争，中国厂商在WAIC 2026上集中展示千卡级SuperPod方案，但软件层AI货币化仍处于早期试点阶段，企业级AI采购周期长于预期，导致头部平台与中小型公司估值分化加剧。}
\subsubsection*{主流共识}
\begin{itemize}
  \item AI系统级互联（scale-up）成为竞争焦点，中国厂商自研NVLink类技术，华为Atlas 950实现1024 NPU互联与统一内存语义。
  \item 日本半导体设备需求从前端设备主导，覆盖HBM、先进逻辑、大宗DRAM和NAND，SEAJ上调FY3/27销售预期至6.55万亿日元（+26\%）。
  \item 企业级AI变现仍处早期，Workday和Adobe的AI产品尚未对冲核心业务增速放缓，采购周期受合规与审计要求拉长。
  \item ServiceNow的AI ACV已超10亿美元，L1 ITSM自动解决80-85\%服务请求，成为企业AI基础设施化的关键节点。
\end{itemize}
\subsubsection*{分歧与条件差异}
\begin{itemize}
  \item 中国AI芯片厂商在系统层面具备竞争力（光学互联、机架设计），但单芯片受限于晶圆制程，与英伟达NVLink生态的差距仍需观察。
  \item 东南亚电信AI代币业务被JPM视为不可持续的利润来源（易复制、收益流向消费者），而自建AI算力的ST和ISAT可能受益。
  \item Shopify的AI成本（LLM）预计影响2026年订阅毛利率约160bps，但管理层可通过模型组合与商家变现应对，市场对毛利率路径存在分歧。
  \item Adobe免费增值模式拉低ARR可见度，有机ARR增长下调约5亿美元，而ServiceNow的AI ACV环比增长40\%+，显示不同产品周期的分化。
\end{itemize}
\subsubsection*{机构观点}
\paragraph{MS} 2026年WAIC显示中国AI计算竞争焦点从单芯片转向系统级互联与集群可靠性。华为Atlas 950 SuperPod实现1024颗NPU互联，FP8算力约1 EFLOPS，HBM内存池化至256TB。几乎所有国产厂商展示64/128卡SuperPod，依赖自研scale-up技术（类似NVLink）。尽管芯片受限于制程，但在光学互联与机架设计上具备竞争力。
{\small\color{refgray}数据：Atlas 950: 1024 NPU互联, 1 EFLOPS FP8, 96TB HBM池化至256TB\par}
{\small\color{refgray}数据：每个64 NPU柜: 64 PFLOPS FP8, 128 PFLOPS FP4, \textasciitilde{}6TB HBM\par}
{\small\color{refgray}数据：UnifiedBus 2.0实现逻辑扁平化互联\par}
{\small\color{refgray}边际变化：相比2025年WAIC以单芯片发布和推理需求为主，2026年焦点转向SuperPod系统方案，竞争维度从芯片规格升级为系统级互联与内存共享。\par}
\paragraph{MS} 蔚来旗下GeniTech（神玑）从智能驾驶芯片供应商扩展为覆盖智能驾驶、具身智能和智能体推理的全域AI芯片平台。NX9031X已出货超30万颗，中端NX9031U提供800 TOPS（风冷）。2025年底开始向第三方授权技术，外部融资约30亿元，估值83亿元。自研芯片替代进口高端计算芯片，每生产一片摊销固定研发成本。
{\small\color{refgray}数据：NX9031X累计出货超30万颗\par}
{\small\color{refgray}数据：NX9031U: 5nm车规制程, 800 TOPS（风冷）\par}
{\small\color{refgray}数据：外部融资约30亿元, 估值83亿元, 蔚来持股63\%\par}
{\small\color{refgray}数据：一颗NX9031算力相当于4颗英伟达Orin\par}
{\small\color{refgray}边际变化：GeniTech从蔚来内部R\&D成本中心转变为可对外创收的AI芯片平台，商业模式从自用扩展至技术授权，可寻址市场从车规芯片拓宽至机器人、智能体推理。\par}
\paragraph{MS} Workday护城河深（年处理超1万亿笔交易，97\%毛利率），但AI货币化仍处早期试点，无法对冲核心HCM增速放缓（FY26降至13\%，预计FY28降至9\%）。管理层放弃中期增长与利润率框架，短期缺乏锚点。AI产品在HR/财务领域采购周期长，客户担忧概率输出与法律责任。估值偏贵（CY27 GAAP P/E 23倍，高于ADSK/CRM/SAP）。
{\small\color{refgray}数据：HCM增速: FY26 13\%, 预计FY28降至9\%\par}
{\small\color{refgray}数据：订阅收入CAGR FY28前仅12-13\%\par}
{\small\color{refgray}数据：CY27 GAAP P/E 23倍, 高于ADSK 20倍/CRM 18倍/SAP 16倍\par}
{\small\color{refgray}边际变化：管理层正式放弃此前分析师日设定的增长与利润率框架，新框架待下次分析师日公布，短期市场失去方向锚点。\par}
\paragraph{GS} 日本半导体设备景气度超预期：4-6月月均销售额5136亿日元（同比+27\%，环比+7\%），前端设备取代后端成为增长主力，需求从HBM/先进逻辑扩展至大宗DRAM和NAND。SEAJ上调FY3/27预测至6.55万亿日元（+26\%），FY3/28至7.40万亿日元（+13\%）。客户催货加速交付趋势持续。GS将美元/日元假设从155调至160，上调5家公司盈利预测。
{\small\color{refgray}数据：4-6月月均销售额5136亿日元, 同比+27\%, 环比+7\%\par}
{\small\color{refgray}数据：SEAJ FY3/27预测: 6.55万亿日元 (+26\%)\par}
{\small\color{refgray}数据：SEAJ FY3/28预测: 7.40万亿日元 (+13\%)\par}
{\small\color{refgray}数据：美元/日元假设从155调至160\par}
{\small\color{refgray}边际变化：前端设备需求从先进制程专用转向大宗存储+先进逻辑双轮驱动，周期敏感型产品线（大宗DRAM、NAND）开始贡献增量，景气度持续性增强。\par}
\paragraph{GS} ServiceNow AI商业化加速：Q2订阅收入同比+23\%（超指引175bps），cRPO同比+21.5\%（超指引200bps）。AI ACV已超10亿美元，重申年底15亿美元目标；AI新ACV环比增长40\%+，拥有AI生产部署的客户数9个月增长9倍。L1 ITSM产品（5月商用）可自动解决80-85\%服务请求，某航空公司迁移全部客服电话至平台（年500万通）。
{\small\color{refgray}数据：Q2订阅收入同比+23\%, 超指引175bps\par}
{\small\color{refgray}数据：cRPO同比+21.5\%, 超指引200bps\par}
{\small\color{refgray}数据：AI ACV超10亿美元, 重申年底15亿美元目标\par}
{\small\color{refgray}数据：AI新ACV环比增长40\%+\par}
{\small\color{refgray}边际变化：AI产品从概念验证进入规模化收入阶段，L1 ITSM和语音AI扩展了应用场景，AI ACV比核心工作流收入更具粘性（加深客户嵌入度）。\par}
\paragraph{JPM} 东南亚电信AI代币业务难以成为持续利润来源：成功产品会被快速复制，收益流向消费者。AIS捆绑10个AI模型定价220泰铢/月（仅为单独订阅成本的6\%），StarHub提供无限AI访问（13.9新元/月起）。真正有价值的是自建AI基础设施（ST计划FY27前投入6亿新元部署11MW AI算力）。台湾远传电信表现突出：1H26移动收入+6\%，ICT+33\%，净利润+17\%（远超全年3\%指引）。
{\small\color{refgray}数据：AIS AI捆绑定价220泰铢/月, 仅为单独订阅成本的6\%\par}
{\small\color{refgray}数据：ST计划FY27前投入6亿新元部署11MW AI算力\par}
{\small\color{refgray}数据：远传电信1H26净利润+17\%, 远超全年3\%指引\par}
{\small\color{refgray}数据：Indosat光纤资产出售: 6.5亿美元, \textasciitilde{}7倍EV/EBITDA\par}
{\small\color{refgray}边际变化：市场对电信AI代币的乐观情绪与JPM的谨慎判断形成对比，后者认为AI代币是用户粘性工具而非利润中心，自建算力才是增量收入来源。\par}
\paragraph{JPM} ASMPT新CEO（来自SkyWater/Intel）上任，聚焦先进封装与TCB业务。TCB迎来多个催化剂：CoWoS/HBM扩产、台积电/OSAT的C2S/C2W订单、全球IDM数据中心CPU订单、HBM4加速爬坡。无焊剂TCB技术领先，预计在韩国客户中拿下30-35\%份额。中国先进封装产能扩张（长电、盛合晶微）和光子学增长提供额外支撑。目标价225港元（30倍远期PE）。
{\small\color{refgray}数据：TCB在韩国客户预计份额30-35\%\par}
{\small\color{refgray}数据：目标价225港元, 基于30倍远期PE\par}
{\small\color{refgray}数据：HBM4E堆叠高度迁移放缓, 12hi仍是主流\par}
{\small\color{refgray}边际变化：新CEO的晶圆代工与先进封装背景暗示ASMPT进一步向IDM/代工厂生态靠拢，SMT业务剥离预期增强，资本重新聚焦先进封装。\par}
\paragraph{JPM} 日经指数反弹至66,000点，但AI半导体板块重新领涨仍需时间。情绪指标降至接近零，但TOPIX未来12个月EPS未明显波动，悲观情绪扩散需新触发因素。CTA止损卖压最严峻阶段可能已过（日经未跌破64,600则风险有限），但高贝塔/高动量策略反弹力度弱，回升集中在少数个股，市场广度不足。市场等待全球AI龙头财报。
{\small\color{refgray}数据：日经指数反弹至66,000点\par}
{\small\color{refgray}数据：情绪指标降至接近零\par}
{\small\color{refgray}数据：CTA多头仓位自5月以来大部分已平仓\par}
{\small\color{refgray}数据：日经64,600为关键支撑位\par}
{\small\color{refgray}边际变化：此轮回调由AI叙事转变和动量交易者调仓驱动，而非盈利预期下修，市场处于卖压减轻但买盘尚未集结的过渡状态。\par}
\paragraph{JPM} Shopify Q2 GMV预计同比+29\%（高于管理层高20\%区间指引），JPM Chase Card数据验证美国CNP支出同比+11\%（较Q1的10\%加速）。企业级商户收入贡献提升约200bps，增长从中小商户向大型企业迁移。毛利率面临AI/LLM成本压力（预计影响2026年订阅毛利率约160bps），但管理层可通过模型组合与商家变现应对。AI竞争方面，商家不愿让大平台完全管理店铺，Shopify的20年商业数据构成壁垒。Sidekick周活跃店铺同比增长4倍。
{\small\color{refgray}数据：Q2 GMV预计同比+29\%, 高于彭博一致预期28\%\par}
{\small\color{refgray}数据：美国CNP支出Q2同比+11\%, 较Q1的10\%加速\par}
{\small\color{refgray}数据：AI/LLM成本预计影响2026年订阅毛利率约160bps\par}
{\small\color{refgray}数据：企业级商户收入贡献提升约200bps\par}
{\small\color{refgray}边际变化：市场对AI竞争（Meta Business Agent）的担忧与JPM的调研结论形成分歧，后者认为Shopify的商户数据与需求转化飞轮构成差异化优势。\par}
\paragraph{JEF} AI系统正从计算设备转向通信系统，数据移动速度成为瓶颈。共封装光学（CPO）的核心是降低SerDes/I/O功耗，而非替代铜互连。TSMC COUPE技术实现芯片间接口低阻抗、16倍键合密度提升、85\%杂散电容降低，同等功耗下速度提升170\%+或功耗降低40\%。边缘AI不是云AI的补充，而是云AI在功耗/延迟/安全瓶颈下的必然延伸：2025年AI半导体市场约2360亿美元（云1600亿+边缘800亿），2028年预计3710亿（边缘1200亿）。边缘AI更可能采用DDR×3D堆叠而非HBM。
{\small\color{refgray}数据：TSMC COUPE: 16倍键合密度提升, 85\%杂散电容降低\par}
{\small\color{refgray}数据：同等功耗下速度提升170\%+或功耗降低40\%\par}
{\small\color{refgray}数据：2025年AI半导体市场2360亿美元（云1600亿+边缘800亿）\par}
{\small\color{refgray}数据：2028年预计3710亿美元（边缘1200亿）\par}
{\small\color{refgray}边际变化：AI系统瓶颈从计算算力转向数据移动效率，CPO和边缘AI成为新投资主题，HBM与边缘AI技术路线分化（HBM加速AI vs DDR×3D堆叠分发AI）。\par}
\paragraph{MS} Adobe面临三重转型叠加（免费增值模式、CEO/CFO交接、AI再投入），ARR增速从FY25的约11.5\%持续放缓。有机ARR增长下调约5亿美元（免费增值与推迟涨价各占一半）。免费MAU从5000万增至9000万，但流量被导向免费体验而非付费转化。利润率接近峰值（AI投入与低ARPU用户增长压缩空间）。估值约12倍GAAP PE已反映不确定性，但缺乏清晰扭转路径。
{\small\color{refgray}数据：有机ARR增长下调约5亿美元\par}
{\small\color{refgray}数据：免费MAU从5000万增至9000万\par}
{\small\color{refgray}数据：Adobe.com商务/消费者流量同比+35\%\par}
{\small\color{refgray}数据：GAAP PE约12倍\par}
{\small\color{refgray}边际变化：免费增值策略主动拉长付费转化路径，ARR可见度下降；管理层更替在战略转型期增加执行不确定性，市场需等待新团队证明AI变现能力。\par}
\paragraph{MS} 北美互联网板块进入Q2财报季前整体下跌4\%，头部平台估值折价显著：亚马逊26倍PE（低于均值17\%）、谷歌23倍PE（低于均值11\%）、Meta 20倍PE（低于均值16\%）。数字广告板块近一周下跌3.1\%（Reddit -7.2\%，Meta -3.5\%，谷歌-2.9\%），电商板块上涨0.7\%（亚马逊+0.8\%）。中小市值公司估值压缩更明显（Pinterest 39倍PE，Reddit 26倍PE，Snap亏损）。
{\small\color{refgray}数据：互联网板块整体下跌4\%, 标普500跌2\%, 纳指跌4\%\par}
{\small\color{refgray}数据：亚马逊26倍PE, 低于均值17\%\par}
{\small\color{refgray}数据：谷歌23倍PE, 低于均值11\%\par}
{\small\color{refgray}数据：Meta 20倍PE, 低于均值16\%\par}
{\small\color{refgray}边际变化：头部平台估值折价反映市场在利率环境下重新定价远期现金流，而非基本面恶化；板块内部分化（广告承压vs电商韧性）预示Q2业绩分化。\par}
\subsubsection*{关键数据}
\begin{itemize}
  \item 华为Atlas 950: 1024 NPU互联, 1 EFLOPS FP8, 96TB HBM池化至256TB
  \item 蔚来GeniTech NX9031X出货超30万颗, 外部融资30亿元, 估值83亿元
  \item Workday HCM增速FY26降至13\%, 预计FY28降至9\%
  \item 日本半导体设备4-6月月均销售额5136亿日元, 同比+27\%
  \item ServiceNow AI ACV超10亿美元, AI新ACV环比增长40\%+
  \item AIS AI捆绑定价220泰铢/月, 仅为单独订阅成本的6\%
  \item ASMPT TCB在韩国客户预计份额30-35\%, 目标价225港元
  \item Shopify Q2 GMV预计同比+29\%, AI/LLM成本影响毛利率约160bps
  \item TSMC COUPE: 16倍键合密度提升, 85\%杂散电容降低, 速度提升170\%+
  \item Adobe有机ARR下调约5亿美元, 免费MAU从5000万增至9000万
\end{itemize}
\begin{figure}[htbp]
\centering
\includegraphics[width=0.88\linewidth]{figures/fig_001.jpg}
\caption{Exhibit 1 - Exhibit 1: Huawei showcased its new Atlas 950 SuperPod at 2026 WAIC Exhibit 2: NPU blade (tray) of Atlas 950 SuperPod Exhibit 3: The reverse side of Atlas 950 mainly showcases the optical interconnects adopted by UBlink}
\end{figure}
\begin{figure}[htbp]
\centering
\includegraphics[width=0.88\linewidth]{figures/fig_045.jpg}
\caption{Exhibit 1 - Exhibit 1: Japan SPE sales (3-MMA) Exhibit 2: Ratings and 12m target prices for our SPE coverage \textbackslash{}*APAC Conviction List}
\end{figure}
\begin{figure}[htbp]
\centering
\includegraphics[width=0.88\linewidth]{figures/fig_063.jpg}
\caption{Figure 1 - Figure 1: ASEAN \& Taiwan telcos and towercos: Total shareholder return over the last month Figure 2: ASEAN \& Taiwan telcos and towercos: Total shareholder return over the last 12 months Figure 3: ASEAN \& Taiwan telcos and towe}
\end{figure}
\begin{figure}[htbp]
\centering
\includegraphics[width=0.88\linewidth]{figures/fig_093.jpg}
\caption{Exhibit 4 - Exhibit 4: NTM EV/EBITDA for AMZN/GOOGL/META vs. Historical Averages Exhibit 5: Internet names fell -4\% last week (SPX/NDX -2/-4\%) Exhibit 6: NTM EV/EBITDA Multiples Are -6\%/-13\% vs. 5/10-Year Averages...}
\end{figure}
{\small\color{refgray}本节综合 12 篇研究；涉及 MS、GS、JPM、JEF。}
\newpage
\subsection{科技与AI（2）：地平线业绩分化，下半年芯片放量成关键}
\textbf{地平线机器人上半年收入符合预期但核心亏损扩大，市场已定价最差情景；下半年芯片收入需同比增长超80\%才能达标，客户订单突破与HSD 2.0商业化进展是估值修复的核心催化剂。}
\subsubsection*{主流共识}
\begin{itemize}
  \item 上半年收入增速25\%-35\%符合管理层此前指引，市场预期已基本反映。
  \item 毛利率60\%-66\%持平或略低，产品解决方案收入占比提升是主要拖累。
  \item 下半年芯片驱动收入需同比增长80\%以上，高度依赖理想、比亚迪、吉利、奇瑞等客户出货扩张。
\end{itemize}
\subsubsection*{分歧与条件差异}
\begin{itemize}
  \item 部分观点认为当前估值已充分反映最差情景，上行空间取决于客户突破；另一些观点则担忧亏损扩Daiwa研发投入持续高企，短期盈利改善无望。
  \item 对HSD 2.0高阶智驾方案的商业化节奏存在分歧：乐观者认为其将打开新增长曲线，悲观者认为落地时间与订单规模仍不确定。
\end{itemize}
\subsubsection*{机构观点}
\paragraph{MS} MS认为地平线上半年业绩预告显示收入增长25\%-35\%至19.3-20.8亿元，符合预期，但剔除股份支付和公允价值变动后净亏损扩大至14-17亿元，同比增5\%-28\%，主因研发投入集中于J7芯片和HSD方案。毛利率60\%-66\%持平或略低，反映产品解决方案占比提升。当前市场定价已反映最差情景，后续估值修复取决于比亚迪等关键客户订单份额突破及HSD 2.0商业化进展。下半年芯片驱动收入需同比增长超80\%才能达到全年一致预期，客户集中度是关键变量。
{\small\color{refgray}数据：1H26收入19.3-20.8亿元，同比+25\%-35\%\par}
{\small\color{refgray}数据：1H26核心净亏损14-17亿元，同比扩大5\%-28\%\par}
{\small\color{refgray}数据：毛利率60\%-66\%，持平或略低于1H25\par}
{\small\color{refgray}数据：下半年芯片驱动收入需同比+80\%以上\par}
{\small\color{refgray}边际变化：此前市场关注收入增速放缓，本次业绩预告显示亏损扩大幅度超预期，且毛利率承压信号更明确；下半年80\%增速目标成为新焦点。\par}
\subsubsection*{关键数据}
\begin{itemize}
  \item 1H26收入19.3-20.8亿元，同比+25\%-35\%
  \item 1H26核心净亏损14-17亿元，同比扩大5\%-28\%
  \item 毛利率60\%-66\%，持平或略低于1H25
  \item 下半年芯片驱动收入需同比+80\%以上
\end{itemize}
{\small\color{refgray}本节综合 1 篇研究；涉及 MS。}
\newpage
\subsection{消费与零售：结构分化与动量冲击下的行业再定价}
\textbf{全球消费零售板块正经历结构性分化：中国快递物流业从价格战转向质量驱动，盈利拐点已确认；而美国运动服饰龙头Lululemon面临美洲业务规模高估风险，盈利预期下修压力显著；同时，美国散户资金行为从被动配置转向主动交易，动量因子极端回撤加剧了科技消费相关标的的波动。}
\subsubsection*{主流共识}
\begin{itemize}
  \item 中国快递物流行业竞争格局从价格战转向质量战，收入增速连续两个月跑赢件量增速，单价回升至7.77元，头部企业盈利拐点已确认。
  \item 美国散户资金整体参与度处于近一年第35百分位，ETF资金流入降至第4百分位，显示被动配置热情降温。
  \item 动量因子因极端拥挤出现剧烈回撤，过去一个月回撤约14\%，并伴随多次高幅度负向尾部事件。
\end{itemize}
\subsubsection*{分歧与条件差异}
\begin{itemize}
  \item Lululemon美洲业务可持续规模存在分歧：MS通过三种独立框架交叉验证认为其当前78亿美元体量可能高估25\%-60\%，而市场普遍预期仍高于70亿美元。
  \item 中国快递行业内部策略分化：申通、圆通通过份额扩张实现盈利高增，韵达则主动选择保利润、舍份额，市占率微降。
  \item 美国散户资金在科技ETF与个股间出现背离：科技ETF流入转负，半导体ETF遭集中卖出，但散户仍净买入NVDA、TSLA等个股。
\end{itemize}
\subsubsection*{机构观点}
\paragraph{MS} MS认为Lululemon美洲业务约78亿美元的体量可能比可持续收入基础高出25\%-60\%，新管理层或需采取“调整以求增长”策略，引发盈利预期下修。三种独立框架（同行对比、正常化增长、规模分析）均指向美洲业务合理规模在30-65亿美元区间，低于市场2026年预期。预计中期每股盈利10-11美元，低于市场预期的12美元。
{\small\color{refgray}数据：Lululemon美洲业务当前约78亿美元\par}
{\small\color{refgray}数据：同行对比框架指向合理规模约50亿美元\par}
{\small\color{refgray}数据：单一品牌时尚公司对比指向约30亿美元\par}
{\small\color{refgray}数据：正常化增长模型指向约65亿美元\par}
{\small\color{refgray}边际变化：此前市场普遍预期Lululemon美洲业务将温和复苏，MS此次通过多框架交叉验证，明确指出当前规模可能高估25\%-60\%，盈利预期下修路径不确定性上升。\par}
\paragraph{JPM} JPM指出中国快递物流行业6月数据确认收入增速（+7.7\%）持续跑赢件量增速（+3.8\%），单价回升至7.77元，竞争格局从价格战转向质量战。圆通、申通、韵达二季度净利润同比分别增长92\%、152\%和136\%，盈利拐点已可验证。未来3-6个月排序重新洗牌，满帮与极兔进入首选，中通、京东物流与顺丰紧随其后。
{\small\color{refgray}数据：6月行业收入同比增长7.7\%至1360亿元\par}
{\small\color{refgray}数据：6月行业件量同比增长3.8\%至175.1亿件\par}
{\small\color{refgray}数据：单价回升至7.77元，同比提升3.8\%\par}
{\small\color{refgray}数据：申通件量同比增长19\%，市占率提升1.8个百分点至14.8\%\par}
{\small\color{refgray}边际变化：此前行业仍处于价格战预期中，6月数据连续第二个月收入增速跑赢件量，单价回升，头部企业盈利高增，标志着竞争格局实质性切换至质量驱动。\par}
\paragraph{JPM} JPM零售雷达周报显示，截至7月22日当周美国散户资金净流入57亿美元，低于12周均值68亿美元，参与度处于第35百分位。结构上ETF资金流入降至第4百分位，而个股交易活跃度升至第72百分位，散户行为从被动配置转向主动交易。动量因子因极端拥挤在过去一个月回撤约14\%，出现多次高幅度尾部事件，包括一次-6西格玛交易日。科技ETF流入转负，SOXL和SMH遭集中卖出，但散户仍净买入NVDA、TSLA等个股。
{\small\color{refgray}数据：散户资金净流入57亿美元，低于12周均值68亿美元\par}
{\small\color{refgray}数据：ETF资金流入降至第4百分位\par}
{\small\color{refgray}数据：个股交易活跃度升至第72百分位\par}
{\small\color{refgray}数据：动量因子过去一个月回撤约14\%\par}
{\small\color{refgray}边际变化：此前散户资金以被动ETF配置为主，本周ETF流入骤降至第4百分位，个股交易升至第72百分位，显示行为模式从配置转向主动交易，且动量因子回撤幅度和速度均超过2025年2-3月。\par}
\subsubsection*{关键数据}
\begin{itemize}
  \item Lululemon美洲业务当前约78亿美元，MS认为可能高估25\%-60\%
  \item 中国快递6月收入同比+7.7\%，件量+3.8\%，单价回升至7.77元
  \item 圆通/申通/韵达二季度净利润同比+92\%/+152\%/+136\%
  \item 美国散户资金净流入57亿美元，低于12周均值68亿美元
  \item 动量因子一个月回撤14\%，出现-6西格玛单日事件
  \item 科技ETF流入转负，SOXL和SMH遭-1.6西格玛卖出
\end{itemize}
\begin{figure}[htbp]
\centering
\includegraphics[width=0.88\linewidth]{figures/fig_006.jpg}
\caption{Exhibit 1 - Exhibit 3: Our scale analysis indicates that LULU Americas has the characteristics of a business that is typically \textbackslash{}\$3-7B in size – below its current \textbackslash{}\$7B + levels.}
\end{figure}
\begin{figure}[htbp]
\centering
\includegraphics[width=0.88\linewidth]{figures/fig_067.jpg}
\caption{Figure 1 - Figure 1: China logistics share price performance (past one month) Figure 2: China logistics share price performance (year to date)}
\end{figure}
\begin{figure}[htbp]
\centering
\includegraphics[width=0.88\linewidth]{figures/fig_082.jpg}
\caption{Figure 1 - Figure 1: Momentum Crowding Figure 2: Momentum Unwind Now vs Feb-Mar'25 Figure 3: Frequent Negative High-Sigma Events in Momentum}
\end{figure}
\begin{figure}[htbp]
\centering
\includegraphics[width=0.88\linewidth]{figures/fig_007.jpg}
\caption{Exhibit 1 - Exhibit 3: Our scale analysis indicates that LULU Americas has the characteristics of a business that is typically \textbackslash{}\$3-7B in size – below its current \textbackslash{}\$7B + levels.}
\end{figure}
{\small\color{refgray}本节综合 3 篇研究；涉及 MS、JPM。}
\newpage
\subsection{工业与能源：周期复苏与结构性转型的交汇点}
\textbf{全球工业与能源领域正经历周期性与结构性力量的双重驱动：模拟芯片与汽车行业显示复苏信号，电网升级与数据中心冷却需求提供长期增长空间，而中国能源体系通过库存调节与结构转换展现出韧性。}
\subsubsection*{主流共识}
\begin{itemize}
  \item 模拟芯片行业周期性复苏正在加速，工业与汽车终端市场同步改善。
  \item 全球电网投资结构性增长，老旧基础设施替换与可再生能源并网是核心驱动力。
  \item 中国能源体系通过库存释放和结构转换有效缓冲了进口下降的影响。
\end{itemize}
\subsubsection*{分歧与条件差异}
\begin{itemize}
  \item 特斯拉资本开支周期与自由现金流压力：市场对持续投入的容忍度取决于Robotaxi与Optimus的里程碑验证。
  \item 模拟芯片毛利率恢复节奏分化：德州仪器因库存偏高恢复较慢，而微芯、恩智浦等同行可能更快。
  \item 中国汽车市场订单回暖强度：新车型驱动个别品牌爆发，但全行业全面复苏尚未确认。
  \item 3M增长可持续性：新产品管线密集交付能否支撑中期增长加速，市场对利润率改善幅度存在分歧。
\end{itemize}
\subsubsection*{机构观点}
\paragraph{MS} 特斯拉正经历超预期的资本开支扩张，2027年资本开支指引接近300亿美元，覆盖Robotaxi、Optimus、半导体工厂和AI算力，但自由现金流消耗预计达140亿美元。FSD北美选装率达55\%，远超预期的25-30\%，但Robotaxi商业化仍需密度与安全数据验证。Optimus即将量产，但供应链复杂度影响爬坡速度。长期价值取决于可验证的里程碑。
{\small\color{refgray}数据：2027年资本开支指引接近300亿美元\par}
{\small\color{refgray}数据：2027年自由现金流消耗预计约140亿美元\par}
{\small\color{refgray}数据：FSD北美选装率55\%，远超预期的25-30\%\par}
{\small\color{refgray}数据：无人监督Robotaxi累计行驶38万英里\par}
{\small\color{refgray}边际变化：资本开支指引大幅上调，市场关注点从成本控制转向ROI实现路径。\par}
\paragraph{MS} 通用汽车二季度业绩超预期，上调2026年EBIT指引中值5亿美元至140-160亿美元。软件递延收入年底将达75亿美元，同比增长50\%，2027年实现两位数增长，毛利率约70\%。国防业务收入接近7亿美元。管理层首次给出2027年方向性展望，收入、EBIT、自由现金流均有望增长。
{\small\color{refgray}数据：2026年EBIT指引中值上调5亿美元至140-160亿美元\par}
{\small\color{refgray}数据：软件递延收入年底达75亿美元，同比+50\%\par}
{\small\color{refgray}数据：2027年软件收入两位数增长，毛利率约70\%\par}
{\small\color{refgray}数据：国防业务收入接近7亿美元\par}
{\small\color{refgray}边际变化：市场关注点从传统汽车制造转向软件服务和非汽车业务，软件业务有望推动估值重估。\par}
\paragraph{MS} 7月第三周中国汽车市场订单回暖，但非全面启动。小鹏和理想凭借新车型实现周订单环比增长585\%和225\%，比亚迪订单企稳（周订单7.67-7.72万，同比+7\%），零跑持续增长（同比+58\%）。鸿蒙智行环比下滑。新车型驱动订单爆发，但可持续性待观察。
{\small\color{refgray}数据：小鹏周订单4.18-4.2万，环比+585\%\par}
{\small\color{refgray}数据：理想周订单1.79-1.81万，环比+225\%\par}
{\small\color{refgray}数据：比亚迪周订单7.67-7.72万，同比+7\%\par}
{\small\color{refgray}数据：零跑周订单1.61-1.66万，同比+58\%\par}
{\small\color{refgray}边际变化：订单回暖由个别新车型驱动，而非全行业复苏，可持续性是关键。\par}
\paragraph{UBS} 3M二季度有机销售额同比增长5.4\%，超出市场预期150个基点，每股收益2.40美元，较预期高7\%。公司计划2026年推出超过350款新产品，到2027年累计超1000款。UBS上调2027年每股收益预期至10.46美元，比市场共识高约10\%，基于5\%有机增长、35\%核心边际利润率和6亿美元净效率提升。
{\small\color{refgray}数据：二季度有机销售额同比+5.4\%，超预期150个基点\par}
{\small\color{refgray}数据：每股收益2.40美元，较预期高7\%\par}
{\small\color{refgray}数据：2027年每股收益预期10.46美元，比共识高10\%\par}
{\small\color{refgray}数据：2026年计划推出350+款新产品\par}
{\small\color{refgray}边际变化：管理层上调全年指引，新产品管线进入密集交付期，增长加速信号明确。\par}
\paragraph{GS} 华明装备报价自高点回调52\%，GS上调评级。全球电网投资2024-2035年复合增速约7\%，公司作为OLTC双寡头之一，全球份额预计从13\%升至16\%。海外业务毛利率比国内高5-10个百分点，出口占比从41\%提升至59\%，推动净利润2026-2030年复合增速约17\%。
{\small\color{refgray}数据：报价自高点回调52\%\par}
{\small\color{refgray}数据：全球电网投资2024-2035年复合增速约7\%\par}
{\small\color{refgray}数据：全球OLTC份额预计从13\%升至16\%\par}
{\small\color{refgray}数据：海外业务毛利率比国内高5-10个百分点\par}
{\small\color{refgray}边际变化：市场将华明与AI基础设施主题捆绑定价，但实际受益于更广泛的电网升级周期。\par}
\paragraph{GS} 2026年二季度中国原油净进口同比下降24\%，天然气和煤炭净进口分别下降7\%和24\%，但总能源需求仍同比正增长0.4\%。缓冲机制包括：煤炭有效去库存贡献3.0个百分点，原油去库存贡献2.2个百分点，电动车渗透率提升使汽油消费下降23\%而充电量增长60\%。
{\small\color{refgray}数据：原油净进口同比-24\%\par}
{\small\color{refgray}数据：煤炭净进口同比-24\%\par}
{\small\color{refgray}数据：总能源需求同比+0.4\%\par}
{\small\color{refgray}数据：煤炭去库存贡献3.0个百分点\par}
{\small\color{refgray}边际变化：中国从全球能源价格冲击的被动接受者转变为主动减震器，库存调节和结构转换能力超预期。\par}
\paragraph{GS} 德州仪器二季度营收54.6亿美元，毛利率61.4\%，每股收益2.14美元，均超预期。工业收入环比增长约10\%，汽车环比高个位数增长，数据中心环比增长20\%。三季度营收指引中值59.0亿美元，高于预期。模拟芯片周期性复苏加速，但毛利率恢复节奏将分化。
{\small\color{refgray}数据：二季度营收54.6亿美元，超预期52.6亿美元\par}
{\small\color{refgray}数据：毛利率61.4\%，超预期59.5\%\par}
{\small\color{refgray}数据：工业收入环比+10\%\par}
{\small\color{refgray}数据：数据中心收入环比+20\%\par}
{\small\color{refgray}边际变化：工业与汽车终端市场首次同步改善，去库存周期基本结束，但毛利率恢复速度因公司而异。\par}
\paragraph{HSBC} 汉钟精机半导体真空泵收入预计从2025年约1亿元增至2026年2.02亿元，2027年增速可能更快。数据中心磁悬浮压缩机市场2028年预计达169亿元，公司正与Trane、Carrier、Vertiv等全球龙头接触测试。工业热泵长期逻辑不变，欧洲市场是未充分定价的增量方向。
{\small\color{refgray}数据：半导体真空泵收入2025年约1亿元，2026年预计2.02亿元\par}
{\small\color{refgray}数据：数据中心磁悬浮压缩机市场2028年预计169亿元\par}
{\small\color{refgray}数据：2026年PE约24倍，2027年降至17.6倍\par}
{\small\color{refgray}边际变化：市场将汉钟视为半导体标的，但实际增长逻辑横跨数据中心冷却、工业热泵等多个大市场。\par}
\subsubsection*{关键数据}
\begin{itemize}
  \item 特斯拉2027年资本开支指引接近300亿美元，自由现金流消耗约140亿美元
  \item FSD北美选装率55\%，远超预期的25-30\%
  \item 3M二季度有机销售额同比+5.4\%，超预期150个基点
  \item 全球电网投资2024-2035年复合增速约7\%
  \item 中国原油净进口同比-24\%，但总能源需求同比+0.4\%
  \item 德州仪器工业收入环比+10\%，数据中心收入环比+20\%
  \item 通用汽车软件递延收入年底达75亿美元，同比+50\%
  \item 数据中心磁悬浮压缩机市场2028年预计169亿元
\end{itemize}
\begin{figure}[htbp]
\centering
\includegraphics[width=0.88\linewidth]{figures/fig_013.jpg}
\caption{Exhibit 2 - Exhibit 2: TSLA Robotaxi Fleet Size ('000)}
\end{figure}
\begin{figure}[htbp]
\centering
\includegraphics[width=0.88\linewidth]{figures/fig_025.jpg}
\caption{Exhibit 1 - Exhibit 1: Huaming is trading at 18.1x 12-m forward P/E, slightly below its historical average of 20.3x since 2015 Exhibit 2: Huaming has pulled back \$52\textbackslash{}\%\$ from the peak, correcting more than its peers, Sieyuan and Yingliu Exh}
\end{figure}
\begin{figure}[htbp]
\centering
\includegraphics[width=0.88\linewidth]{figures/fig_040.jpg}
\caption{Exhibit 1 - Exhibit 1: Coal and Oil Effective Destocking and Higher Renewables Usage Helped Buffer China's Negative Fossil Fuel Import Shock The units here are percentage point contributions to total energy demand growth. “Oil Supply” is dec}
\end{figure}
\begin{figure}[htbp]
\centering
\includegraphics[width=0.88\linewidth]{figures/fig_087.jpg}
\caption{Exhibit 1 - Exhibit 1: 2027 MS EBIT Bridge Software and Services is an underappreciated growth driver. GM is positioning the digital services vector as an increasingly material and less-cyclical earnings engine, with 1mn new digital subscrip}
\end{figure}
{\small\color{refgray}本节综合 8 篇研究；涉及 MS、UBS、GS、HSBC。}
\newpage
\subsection{医疗与生物科技：资金轮动与数据催化下的结构性机会}
\textbf{中国医疗健康板块近期受益于AI主题资金外溢与CXO板块结构性增持，同时康方生物HARMONi研究OS数据改善提升了海外获批概率，两者共同构成短期催化剂，但政策不确定性与资金持续性仍是市场分歧焦点。}
\subsubsection*{主流共识}
\begin{itemize}
  \item 中国医疗健康板块自6月26日低点反弹，A股涨11\%、H股涨16\%，跑赢恒生指数7个百分点、沪深300指数19个百分点，主要由AI主题资金流出后的轮动配置驱动。
  \item CXO板块成为本轮资金迁移的主要受益者，北向资金二季度持仓占比从7.03\%升至9.93\%，环比增加2.9个百分点，而制药和医疗器械板块持仓下降。
  \item 康方生物HARMONi研究OS数据持续改善，全ITT人群HR从0.79降至0.76，西方患者亚组HR从0.98大幅降至0.76，与国内HARMONi-A研究（HR 0.74）基本可比，增加了获批概率。
\end{itemize}
\subsubsection*{分歧与条件差异}
\begin{itemize}
  \item 资金轮动持续性存疑：若AI重新成为主线，医疗健康板块可能面临资金再次流出；部分机构认为当前回升属于价值发现，滞后于上半年强劲的BD势头。
  \item 国内政策不确定性：生物类似药集采、耗材集采、医疗服务费调整等可能压制板块估值修复空间。
  \item 美国《生物安全法案》对跨境BD交易的影响尚未完全消化，可能制约CXO板块的长期增长叙事。
\end{itemize}
\subsubsection*{机构观点}
\paragraph{NOM} 康方生物HARMONi研究OS数据持续改善，全ITT人群HR从2025年4月的0.79降至2026年6月的0.76，西方患者亚组HR从0.98大幅降至0.76，中位随访时间延长至23.2个月，数据成熟度显著提升。NOM认为这一变化增加了该适应症（2L+ EGFRm NSCLC after 3rd gen TKI）获批的可能性，PDUFA日期为2026年11月14日。基于DCF模型维持目标价142.50港元（WACC 10.8\%，终端增长率4\%）。
{\small\color{refgray}数据：全ITT人群OS HR从0.79（2025年4月）降至0.76（2026年6月）\par}
{\small\color{refgray}数据：西方患者亚组OS HR从0.98（2025年4月）降至0.76（2026年6月）\par}
{\small\color{refgray}数据：西方患者亚组中位随访时间从9.2个月（2025年4月）延长至23.2个月（2026年6月）\par}
{\small\color{refgray}数据：HARMONi-A研究OS HR为0.74，与最新HARMONi数据基本可比\par}
{\small\color{refgray}边际变化：OS数据成熟度提升，西方患者亚组HR从0.98大幅改善至0.76，与全ITT人群一致，显著提高了获批概率预期。\par}
\paragraph{HSBC} 中国医疗健康板块自6月26日低点反弹，A股涨11\%、H股涨16\%，跑赢恒生指数7个百分点、沪深300指数19个百分点。上涨主要由AI主题资金流出后的轮动配置驱动，CXO板块成为北向资金主要受益者，二季度持仓占比从7.03\%升至9.93\%，环比增加2.9个百分点。全球投资者最关注CDMO和生物科技/制药领域，药明康德、药明生物、康龙化成因上半年业绩预期较高及2026-2028年增长加速而受青睐。
{\small\color{refgray}数据：A股医疗健康指数自6月26日低点反弹11\%，H股反弹16\%\par}
{\small\color{refgray}数据：跑赢恒生指数7个百分点，跑赢沪深300指数19个百分点\par}
{\small\color{refgray}数据：北向资金CXO持仓占比从一季度的7.03\%升至二季度的9.93\%\par}
{\small\color{refgray}数据：制药和医疗器械板块北向持仓分别下降0.78和0.45个百分点\par}
{\small\color{refgray}边际变化：资金从AI主题向非AI板块轮动，医疗健康成为集中讨论方向之一，CXO板块获得结构性增持，而此前市场对医疗板块关注度较低。\par}
\subsubsection*{关键数据}
\begin{itemize}
  \item 全ITT人群OS HR从0.79（2025年4月）降至0.76（2026年6月）
  \item 西方患者亚组OS HR从0.98（2025年4月）降至0.76（2026年6月）
  \item 西方患者亚组中位随访时间从9.2个月延长至23.2个月
  \item A股医疗健康指数自6月26日低点反弹11\%，H股反弹16\%
  \item 跑赢恒生指数7个百分点，跑赢沪深300指数19个百分点
  \item 北向资金CXO持仓占比从7.03\%升至9.93\%
  \item 制药和医疗器械板块北向持仓分别下降0.78和0.45个百分点
\end{itemize}
\begin{figure}[htbp]
\centering
\includegraphics[width=0.88\linewidth]{figures/fig_052.jpg}
\caption{Exhibit 1 - Issuer of report: HSBC Qianhai Securities Limited View HSBC Qianhai Securities at: https://www.research.hsbc.com \#\# Industry update Exhibit 1. Y-t-d performance of China Healthcare index Exhibit 2. North Bound Holding and changes}
\end{figure}
{\small\color{refgray}本节综合 2 篇研究；涉及 NOM、HSBC。}
\newpage
\subsection{行业出清与需求结构重塑：中国药房、光纤与全球电网设备}
\textbf{中国连锁药房行业正经历客流驱动的同店增长加速，源于中小药店出清带来的结构性红利；全球光纤市场因AI数据中心需求爆发而价格创历史新高，但中国产能快速扩张预示2026年或为周期顶点；电网设备受益于数据中心订单陡峭增长，订单出货比达1.7倍，积压订单同比增65\%。}
\subsubsection*{主流共识}
\begin{itemize}
  \item 中国连锁药房行业同店增长在2026年第二季度重新加速，驱动力是客流增长而非客单价提升。
  \item 全球光纤市场因AI数据中心需求爆发，价格已突破历史高点，2026年二季度均价超100元/芯公里。
  \item 电网设备需求强劲，数据中心成为新增量引擎，GE Vernova上半年数据中心订单已超2025年全年。
\end{itemize}
\subsubsection*{分歧与条件差异}
\begin{itemize}
  \item 药房行业：市场担忧线上比价和飞行检查等政策风险，但JPM认为这些已从基本面问题转为情绪扰动，行业出清带来的客流再分配才是更实质的驱动。
  \item 光纤行业：JEF认为2026年是盈利顶点，因中国产能快速扩张（2028年全球产能增69\%），而GS对日本银行板块的积极看法则基于贷款增速加速和利率正常化。
  \item 特斯拉：GS认为短期利润率低于预期（毛利率16.8\% vs 共识19.5\%）并非核心，应关注FSD和robotaxi进展；但市场对盈利路径分歧较大。
\end{itemize}
\subsubsection*{机构观点}
\paragraph{JPM} 中国连锁药房行业正经历结构性转折：同店增长在2026年第二季度重新加速至5-6\%，驱动力是客流增长而非客单价。核心逻辑是行业出清——全国药店总数从近70万家预计3-5年缩减至40-50万家，单店服务人口从约2000人提升至约3000人，每年净关店3-5万家可带来约5\%的年化同店增长。并购估值已压缩至0.3-0.5倍市销率，加盟模式取代自建成为扩张主力，新签加盟店中70-80\%为存量门店转化。非药品类（功能性食品、个人护理）成为新增长方向。
{\small\color{refgray}数据：益丰二季度同店增速约6\%，老百姓5-6\%，大参林环比加速\par}
{\small\color{refgray}数据：全国药店总数从近70万家预计缩减至40-50万家\par}
{\small\color{refgray}数据：并购估值压缩至0.3-0.5倍市销率\par}
{\small\color{refgray}数据：新签加盟店中70-80\%为存量门店转化（两年前仅30\%）\par}
{\small\color{refgray}边际变化：市场此前关注政策风险，但行业出清带来的客流再分配成为更实质的驱动因素，同店增长从低个位数加速至5-6\%。\par}
\paragraph{JEF} 全球光纤市场因AI数据中心需求爆发而价格创历史新高，2026年二季度均价已突破100元/芯公里，长飞光纤2026年均价预计达146元/芯公里，较2025年增长203\%，比2018年峰值高18\%。但中国产能快速扩张（2028年全球产能增69\%），2026年或为周期顶点。数据中心占光纤需求比例将从2024年的10\%跃升至2027年的55\%。长飞光纤是目前较少未宣布扩产的主要厂商，但新进入者（大族激光、合盛硅业）带来供给压力。
{\small\color{refgray}数据：光纤均价从2020年31元/芯公里升至2025年44元，2026年二季度突破100元\par}
{\small\color{refgray}数据：长飞光纤2026年均价预计146元/芯公里，较2025年增长203\%\par}
{\small\color{refgray}数据：全球光纤产能到2028年将增长69\%\par}
{\small\color{refgray}数据：数据中心占光纤需求比例从2024年10\%跃升至2027年55\%\par}
{\small\color{refgray}边际变化：光纤价格突破历史高点，驱动因素从电信需求切换为数据中心需求，价格弹性更大；但中国产能快速扩张使2026年成为盈利顶点。\par}
\paragraph{GS} GE Vernova电气化业务2026财年第二季度订单63.47亿美元，同比增长93\%，远超预期。数据中心需求成为新增长引擎，上半年数据中心订单超50亿美元，已超2025财年全年。订单出货比1.7倍，积压订单446亿美元，同比增65\%。EBITDA利润率18.4\%，同比提升390个基点。固态变压器样品下半年向超大规模客户发货，可能带动800VDC相关产品需求。这对业务结构相似的日立能源构成积极信号。
{\small\color{refgray}数据：GE Vernova电气化业务Q2订单63.47亿美元，同比增93\%\par}
{\small\color{refgray}数据：上半年数据中心订单超50亿美元，超2025财年全年\par}
{\small\color{refgray}数据：订单出货比1.7倍，积压订单446亿美元，同比增65\%\par}
{\small\color{refgray}数据：EBITDA利润率18.4\%，同比提升390个基点\par}
{\small\color{refgray}边际变化：数据中心订单从线性增长变为陡峭化，半年超去年全年；固态变压器等新产品开启增量市场。\par}
\paragraph{GS} 特斯拉二季度非GAAP每股收益0.33美元，远低于共识0.55美元和GS预期0.58美元。毛利率16.8\%，低于共识19.5\%。汽车业务毛利率16.3\%，受利率促销和商品价格波动影响。能源业务毛利率20.4\%，剔除2.4亿美元保修费用后约28\%。FSD活跃订阅用户150万，环比增16\%，同比增56\%；北美新车FSD选装率55\%。robotaxi累计行驶38万英里无重大事故。短期利润低于预期，但自动驾驶和机器人进展是更关键的中期驱动。
{\small\color{refgray}数据：Q2非GAAP EPS 0.33美元，低于共识0.55美元\par}
{\small\color{refgray}数据：毛利率16.8\%，低于共识19.5\%\par}
{\small\color{refgray}数据：FSD活跃订阅用户150万，同比增56\%\par}
{\small\color{refgray}数据：北美新车FSD选装率55\%\par}
{\small\color{refgray}边际变化：市场在7月初强劲交付报告后对利润预期偏高，但实际毛利率低于预期；FSD和robotaxi商业化出现可量化进展。\par}
\paragraph{GS} 日本银行板块受五大顺风因素支撑：利率上升（2025年12月加息后抵押贷款利率上调、2026年6月加息前短端利率上行）、贷款增速加快（城市银行从1-3月5.9\%加速至4-6月8.4\%）、活跃资本市场、TOPIX指数上涨带来信托赎回收益、外汇影响。中期ROE前景使P/B看向2倍或更高。读者关注焦点从金融科技回归传统银行，对日本交易所集团（JPX）兴趣上升。
{\small\color{refgray}数据：城市银行贷款增速从1-3月5.9\%加速至4-6月8.4\%\par}
{\small\color{refgray}数据：2025年12月加息后抵押贷款利率上调\par}
{\small\color{refgray}数据：2026年6月加息前短端利率上行\par}
{\small\color{refgray}数据：P/B看向2倍或更高\par}
{\small\color{refgray}边际变化：贷款增速加速成为可跟踪的先行指标；读者关注点从金融科技回归传统银行。\par}
\paragraph{GS} 2026年6月中国广义财政赤字（AFD）继续收窄，3个月移动平均占GDP比重从5月-8.4\%收窄至-7.0\%，12个月移动平均从-10.6\%收窄至-10.2\%。财政政策在二季度实际处于紧缩状态，负向财政脉冲贡献了当季实际GDP环比降幅的40\%以上。预算内收入增速8.7\%高于支出增速4.0\%。土地出让收入同比降幅从5月35.8\%扩大至42.1\%，房地产相关税收降幅从2.6\%扩大至12.3\%。
{\small\color{refgray}数据：6月广义财政赤字3个月移动平均占GDP比重从-8.4\%收窄至-7.0\%\par}
{\small\color{refgray}数据：负向财政脉冲贡献Q2 GDP环比降幅的40\%以上\par}
{\small\color{refgray}数据：6月土地出让收入同比-42.1\%（5月-35.8\%）\par}
{\small\color{refgray}数据：预算内收入增速8.7\%，支出增速4.0\%\par}
{\small\color{refgray}边际变化：广义财政赤字收窄并非主动紧缩，而是土地收入塌陷和准财政工具调整的被动结果；财政脉冲转负对GDP拖累显著。\par}
\paragraph{GS} 耐克将从2027年1月1日起终止两家主要经销商（滔搏、宝胜）在中国大陆的线上平台销售授权。受影响业务分别占滔搏2026财年营收约22\%、宝胜2025财年营收约15\%。但利润冲击幅度差异显著：滔搏线上利润率高于集团整体，利润下调8-35\%；宝胜线上业务稀释利润，剥离后综合毛利率反而可能提升，利润影响7-9\%。耐克渠道调整为安踏、李宁创造份额机会，但需关注耐克清理库存可能引发全行业促销。
{\small\color{refgray}数据：滔搏受影响营收约22\%，宝胜约15\%\par}
{\small\color{refgray}数据：滔搏利润下调8-35\%，宝胜利润下调7-9\%\par}
{\small\color{refgray}数据：耐克中国运动鞋服市场份额约16\%\par}
{\small\color{refgray}数据：滔搏股价已下跌44\%，股息率约13\%\par}
{\small\color{refgray}边际变化：市场此前低估了线上业务利润率差异对利润冲击的影响；耐克渠道调整为本土品牌创造窗口，但促销风险上升。\par}
\subsubsection*{关键数据}
\begin{itemize}
  \item 中国连锁药房行业预计3-5年从近70万家缩减至40-50万家
  \item 光纤均价2026年二季度突破100元/芯公里，长飞预计146元/芯公里
  \item GE Vernova上半年数据中心订单超50亿美元，超2025财年全年
  \item 特斯拉Q2毛利率16.8\%，低于共识19.5\%
  \item 日本城市银行贷款增速从5.9\%加速至8.4\%
  \item 中国6月广义财政赤字收窄至-7.0\%（3个月移动平均）
  \item 耐克线上代理终止影响滔搏营收22\%、宝胜15\%
\end{itemize}
\begin{figure}[htbp]
\centering
\includegraphics[width=0.88\linewidth]{figures/fig_035.jpg}
\caption{Exhibit 1 - Exhibit 1: Fiscal revenue growth continued to outperform spending growth in June Exhibit 2: Property-related government revenue weakened further in June Exhibit 3: Our augmented fiscal deficit (AFD) metric continued to narrow i}
\end{figure}
\begin{figure}[htbp]
\centering
\includegraphics[width=0.88\linewidth]{figures/fig_037.jpg}
\caption{Exhibit 2 - Exhibit 5: Fiscal impulse turned negative in Q2, while we expect it to turn positive in H2}
\end{figure}
\begin{figure}[htbp]
\centering
\includegraphics[width=0.88\linewidth]{figures/fig_036.jpg}
\caption{Exhibit 1 - Exhibit 4: China's fiscal "spend-through" ratio rose in June on a 12mma basis}
\end{figure}
\begin{figure}[htbp]
\centering
\includegraphics[width=0.88\linewidth]{figures/fig_038.jpg}
\caption{Exhibit 3 - Exhibit 5: Fiscal impulse turned negative in Q2, while we expect it to turn positive in H2 \#\# The China Economics Team}
\end{figure}
{\small\color{refgray}本节综合 7 篇研究；涉及 GS、JPM、JEF。}
\newpage
\section{辅助信号｜战略咨询}
本期战略咨询板块聚焦欧洲化工行业的结构性重置与AI算力商品化两大主题。欧洲化工正经历非周期性的竞争力下滑，行业分化加速，组合平均指标失效；AI算力市场则从固定订阅转向实时定价，释放出每年高达1400亿美元的暗价值，企业用户面临成本敞口显著上升的新格局。
\subsection{欧洲化工的结构性重置：非周期性下滑与行业分化}
\textbf{欧洲化工行业正经历一场由能源成本、全球产能过剩、碳监管和需求结构变化驱动的结构性重置，而非周期性波动。行业已分裂为差异化可防御、不确定性过渡型和结构性暴露三个群体，组合平均指标失去指导意义，战略决策必须下沉至细分领域。}
\begin{itemize}
  \item 欧洲化工装置利用率较2015-2020年均值下降约10个百分点，背后是能源与原料成本高企、全球产能过剩、需求模式转变及监管收紧等多重力量叠加，这些不确定性不会随宏观周期消退。
  \item 2026年Q2中东局势导致全球石油供应路线中断，欧洲化工生产商竞争地位一度改善，但波士顿咨询定性为暂时性好转，核心问题——需求调整、投入成本高、全球产能过剩及亚洲/中东低价竞争——并未解决。
  \item 行业分化加速：差异化可防御领域凭借创新和客户粘性维持利润；不确定性过渡型业务成本劣势和监管暴露持续侵蚀竞争力；结构性暴露的大宗化学品（如聚合物）面临能源密集、易于贸易、难以差异化的严峻挑战。
  \item 四个结构性驱动变量：能源强度、贸易暴露度、终端市场周期性和监管负担。高能源强度产品受制于欧洲成本位置，易于运输的化学品面临更强进口竞争，高碳足迹推高监管成本。
  \item 欧洲天然气价格是美国的4倍，电力是2倍；EU ETS碳价五年上涨约40欧元/吨；中国聚丙烯进口占比十年翻倍——这些边际变化共同构成结构性压力，增量改进计划无法弥合竞争力差距。
\end{itemize}
\begin{figure}[htbp]
\centering
\includegraphics[width=0.88\linewidth]{figures/fig_125.jpg}
\caption{EXHIBIT 1 - Regulation adds further cost pressure. Carbon pricing under the European Union's Emissions Trading System has increased by roughly €40 per ton over the past five years, alongside broader regulatory tightening. Taken together, these pressures create a structura}
\end{figure}
\subsection{AI算力商品化：从固定订阅到实时定价，暗价值释放}
\textbf{企业级AI定价模式正从固定订阅费转向按量计费，下一步将是动态定价——算力价格随供需实时波动。这一转变可能释放每年高达1400亿美元的暗价值，集中在芯片、租赁和token三个环节的优化与套利，以及数据中心运营商的融资成本降低。}
\begin{itemize}
  \item Anthropic（2026年4月）将企业客户改为每人20美元基础费加按token付费的API费率；OpenAI Codex从固定消息定价转为按token计量；GitHub收紧Copilot限制；Windsurf取消固定积分制改为每周硬性限额——边际变化明确指向成本敞口上升。
  \item 下一步是动态token定价：推理价格随底层算力供需波动，类似批发电力市场——实时客服请求在早高峰价格更高，夜间批量摘要任务更便宜。有效运行前提是模型可替代，开发者可在Claude、GPT或开源模型间切换。
  \item BCG估计暗价值高达每年1400亿美元，来源包括：算力芯片、租赁和终端token三个环节的优化与套利，以及数据中心运营商融资成本降低。大型企业用户有望捕获其中320亿至460亿美元。
  \item 当前市场存在三重不透明：GPU现货价格差异大（同一GPU在同一天不同指数价格显著不同）、资本可用性仍高（AI实验室股权融资仍可获得）、管理价格不确定性近期才成为优先事项。
  \item 企业需区分可替代与不可替代场景：深度生态集成（定制芯片、微调模型、数据驻留）创造切换成本，削弱现货市场定价约束；早期行动者——管理成本敞口的企业、获取更便宜融资的数据中心、保障供应的AI实验室——有机会捕获不成比例的价值。
\end{itemize}
\begin{figure}[htbp]
\centering
\includegraphics[width=0.88\linewidth]{figures/fig_126.jpg}
\caption{EXHIBIT 1 - Players that move early will be in the best position to adapt to a rapidly changing market and unlock dark value in the AI compute market. \#\# EXHIBIT 1 AI Market Evolution Can Release up to \textbackslash{}\$140 Billion of Value DARK VALUE (\textbackslash{}\$BILLIONS) Sources: LSEG; GS Globa}
\end{figure}
\begin{figure}[htbp]
\centering
\includegraphics[width=0.88\linewidth]{figures/fig_129.jpg}
\caption{EXHIBIT 5 - Time complexity Daily price volatility of Nvidia H100 SXM \textbackslash{}\$/HOUR PER 1 GPU Quality complexity Average daily price aggregated globally \textbackslash{}\$/HOUR PER 1 GPU Geographic complexity Average daily price of Nvidia H100 SXM \textbackslash{}\$/HOUR PER 1 GPU Sources: AWS; GCP; Vast.ai, }
\end{figure}
\newpage
\section{辅助信号｜智库/国际机构}
本板块整合IMF与世界银行最新研究，聚焦三大主题：IMF透明度改革加速信息披露、全球贸易政策进入高频波动期、以及流离失所儿童发展缺口的异质性。
\subsection{IMF透明度改革与信息披露提速}
\textbf{IMF新版透明度政策指南将出版同意机制从“明确同意”转向“推定同意”，并大幅收窄文件修改窗口，系统性缩短国家文件公开时间，强化其作为全球监督机构的即时信息披露能力。}
\begin{itemize}
  \item 监督类国家文件出版同意机制转为“非反对即同意”，成员国不再需主动同意，沉默或拖延将视为同意，显著降低信息延迟风险。
  \item 新闻稿发布节奏从可协商变为有明确截止日：董事会审议后2个工作日内发布，成员国可申请延期至多7个日历日。
  \item 文件修改窗口收窄至董事会审议后7个日历日，28天后不再受理修改请求，已发布文件原则上不允许修改。
  \item 这一变化意味着成员国更难通过“不表态”延迟不利信息公开，IMF在“可信顾问”与“全球监督者”角色间向后者倾斜。
\end{itemize}
\begin{figure}[htbp]
\centering
\includegraphics[width=0.88\linewidth]{figures/fig_098.jpg}
\caption{Figure 1 - \textbackslash{}- 28 calendar days after the date of the Board's consideration. If the member that requested another time extension to decide on publication has not explicitly objected to publication within 28 calendar days from Board consideration, the member's consent to p}
\end{figure}
\begin{figure}[htbp]
\centering
\includegraphics[width=0.88\linewidth]{figures/fig_099.jpg}
\caption{Figure 2 - 17. In specific circumstances, the Fund will issue factual statements informing the public of the Board's consideration and the members' publication decision. Such statements will be issued if (i) the member has not communicated its publication decision to the}
\end{figure}
\subsection{全球贸易政策进入高频波动期}
\textbf{IMF与WTO联合构建的贸易政策活动指数（TPA）显示，2019年后全球贸易政策使用频率和强度出现结构性跃升，贸易促进与限制措施同步增加，政策工具成为各国应对外部冲击的综合性工具箱。}
\begin{itemize}
  \item TPA指数显示2008-2018年全球贸易政策活动平稳，2019年后持续显著上行，由中美关税升级、疫情和乌克兰战争三重冲击叠加驱动。
  \item 贸易促进类措施与限制性措施同步增加，反映各国同时使用多种政策工具调整的普遍模式，贸易政策不再是简单的开放或封闭二选一。
  \item TPA指数整合关税、非关税措施、贸易救济、补贴、海关程序等，通过动态因子模型提取共同趋势，提供可实时监测的指标。
  \item 方法论框架可推广至产业政策、气候政策、数字监管等多源多维度政策监测，先提取共同趋势，再结合高频数据即时预测。
\end{itemize}
\begin{figure}[htbp]
\centering
\includegraphics[width=0.88\linewidth]{figures/fig_106.jpg}
\caption{Figure 1 - Baseline Results: The Overall Trade Policy Activity Index This section presents the new Trade-Policy Activity (TPA) Index that aims to capture in a systematic fashion the dynamics in overall trade policy activity worldwide over time. The Index plotted in Figur}
\end{figure}
\begin{figure}[htbp]
\centering
\includegraphics[width=0.88\linewidth]{figures/fig_107.jpg}
\caption{Figure 2 - \#\# Understanding the Role of Different Types of Measures The global index captures the common movement in all types of trade policy measures, including measures that facilitate trade, such as more efficient border procedures, and those that could restrict or d}
\end{figure}
\subsection{流离失所儿童发展缺口的异质性与政策启示}
\textbf{世界银行基于哥伦比亚麦德林约1600名儿童的一手数据发现，流离失所原因决定发展缺口的维度：委内瑞拉难民儿童在认知和身体发展上落后最严重，而哥伦比亚国内流离失所儿童的心理健康和社会情感缺口更突出，且这些缺口在控制家庭财富后仍然存在。}
\begin{itemize}
  \item 委内瑞拉难民儿童接受性词汇得分比非流离失所儿童低0.34个标准差，年龄标准化身体质量指数低0.50个标准差，与逃离极端物质匮乏直接相关。
  \item 哥伦比亚国内流离失所儿童抑郁得分低0.25个标准差，社会情感功能低0.21个标准差，创伤症状低0.18个标准差，与武装冲突暴露和代际创伤传递一致。
  \item 两类流离失所儿童在五个发展指标上的系数存在显著差异，表明流离失所不是同质化经历，政策干预必须针对具体发展维度设计。
  \item 研究还发现父母间接受到影响的儿童也表现出类似模式，指向创伤的代际传递渠道，提示干预需覆盖家庭层面。
\end{itemize}
\begin{figure}[htbp]
\centering
\includegraphics[width=0.88\linewidth]{figures/fig_111.jpg}
\caption{Figure 1 - Age at displacement. Figure 1 illustrates the distribution of age at displacement for children in our sample. The figure highlights a fundamental difference between the two displaced populations. Nearly all Venezuelan refugee children experienced displacement }
\end{figure}
\begin{figure}[htbp]
\centering
\includegraphics[width=0.88\linewidth]{figures/fig_112.jpg}
\caption{Figure 2 - Two distinct displacement profiles emerge. Refugee children exhibit steep deficits in cognitive and physical development, scoring 0.54 standard deviations below non-displaced children on the PPVT and approximately 0.56 standard deviations below on BMI, yet the}
\end{figure}
\subsection{AI、技术与生产率}
\textbf{用于判断 AI 投资周期、生产率扩散和产业约束是否正在改变资产定价。}
\begin{itemize}
  \item 南非保护区网络覆盖约11\%的国土面积，但关于其长期宏观环境与生态效益的因果证据长期稀缺。世界银行发展研究组近期发布的工作论文，利用机器学习反事实方法，首次将保护区对家庭收入、鸟类生物多样性和旅游消费者剩余的影响纳入统一框架进行估算。报告的核心发现是：南非的保护区在长达数十年的时间...
\end{itemize}
\begin{figure}[htbp]
\centering
\includegraphics[width=0.88\linewidth]{figures/fig_120.jpg}
\caption{Figure 1 - Figure 1: A map of South Africa with protected areas in green, national parks labelled in black, and major cities labelled in blue. To overcome these limitations, we develop a complementary identification strategy grounded in the}
\end{figure}
\subsection{宏观政策与金融稳定}
\textbf{用于校准利率、通胀、财政约束和金融稳定叙事。}
\begin{itemize}
  \item 所罗门群岛正在经历的气候变化并非遥远威胁，而是已经开始重塑其宏观环境基础结构的缓慢变量。IMF 最新一期国别报告以所罗门群岛为对象，系统评估了升温、海平面上升和海洋生态系统变化对该国宏观环境增长、公共财政和关键出口部门的潜在影响。报告的核心判断是：虽然气候变化带来的不确定性分布不...
\end{itemize}
\begin{figure}[htbp]
\centering
\includegraphics[width=0.88\linewidth]{figures/fig_101.jpg}
\caption{Figure 1 - Drawing on recent scientific data and model-based projections, this section provides an overview of the evolving risks and adaptation challenges posed by a warming climate, changes in marine ecosystems, and sea level rise. 2. Solomon Islands experiences high i}
\end{figure}
\subsection{气候、发展与社会结构}
\textbf{用于观察长期增长、资源约束、社会结构和政策执行质量。}
\begin{itemize}
  \item 1960年代初，埃及实施了一项具有典型意义的公共就业保障政策：所有中学和大学毕业生无需竞争性考试即可获得公共部门终身职位。这项政策持续了近二十年，直到1980年代初因财政不确定性和队列规模膨胀而事实上暂停。世界银行研究团队利用1959/60学年地区级中学供给差异作为政策暴露强度的...
\end{itemize}
\begin{figure}[htbp]
\centering
\includegraphics[width=0.88\linewidth]{figures/fig_115.jpg}
\caption{Figure 1 - Employment in the public sector in 1950s Egypt was subject to a public competition examination. Law 210 of 1951 and Law 113 of 1958 determined the procedure for recruitment for state civil servants and in joint stock companies and public institutions, respecti}
\end{figure}
\section{结语}
后续需验证的关键节点包括：日本GPIF资产配置调整的实际节奏与对日元影响；中国AI芯片厂商系统级竞争力与英伟达NVLink生态的差距；企业级AI采购周期是否缩短及货币化拐点；Lululemon美洲业务规模下修风险；中国光纤产能扩张对全球价格的影响；以及美国《生物安全法案》对跨境BD交易的长期制约。
\newpage
\section{覆盖概览}
投行/券商 40 篇；战略咨询 2 篇；智库/国际机构 6 篇。\par
投行覆盖：GS 12篇 / MS 12篇 / JPM 8篇 / HSBC 3篇 / JEF 2篇 / NOM 2篇 / UBS 1篇\par
\section{Disclaimer}
{\scriptsize\color{refgray}
This community is a paid learning and information-sharing community created for general educational purposes only. The membership fee is charged solely for access to community discussions, educational content, organization, and operational costs. It is not an advisory fee, management fee, performance fee, brokerage fee, or compensation for personalized financial, investment, legal, tax, or professional advice.\par
I participate and share information solely as an individual and for educational discussion among community members. I am not acting as a financial adviser, investment adviser, broker, dealer, portfolio manager, legal adviser, tax adviser, accountant, fiduciary, or any other licensed professional adviser. Nothing shared in this community constitutes, and should not be interpreted as, financial advice, investment advice, legal advice, tax advice, a recommendation, solicitation, offer, endorsement, instruction, or guarantee to buy, sell, hold, short, trade, or invest in any security, cryptocurrency, fund, derivative, strategy, or other financial product.\par
All content shared in this community, including messages, discussions, links, files, charts, examples, market commentary, watchlists, AI-generated or AI-polished summaries, and third-party materials, is general, impersonal, and educational in nature. It does not take into account any individual's financial situation, investment objectives, risk tolerance, investment horizon, liquidity needs, tax status, legal restrictions, or suitability requirements. No content should be relied upon as suitable for any specific person, account, portfolio, or financial decision.\par
Information shared in this community may be incomplete, inaccurate, outdated, speculative, AI-generated, AI-edited, or based on third-party internet sources that have not been independently verified. I make no representation or warranty as to the accuracy, completeness, timeliness, reliability, or suitability of any information shared.\par
Financial markets involve substantial risk, including the possible loss of principal. Past performance, historical data, hypothetical examples, simulated results, backtests, personal opinions, or educational examples do not guarantee future results. Each member is solely responsible for conducting their own research, exercising independent judgment, and making their own decisions. Before making any financial, investment, legal, or tax decision, members should consult qualified and licensed professionals.\par
Participation in this community, payment of any membership fee, or communication with me or other members does not create any adviser-client, fiduciary, professional, contractual, agency, partnership, employment, or other advisory relationship. By joining or participating in this community, you acknowledge and agree that you are solely responsible for your own decisions, actions, transactions, profits, losses, risks, and consequences, and that I shall not be liable for any direct, indirect, incidental, consequential, financial, legal, tax, or other loss or damage arising from or related to any content, discussion, information, or material shared in this community.\par
}
\newpage
\begin{center}
{\Large 更多详情报告kcdesk.com}\par\vspace{0.8cm}
\includegraphics[width=0.58\linewidth]{\detokenize{../../prompts/zsxq_img.jpg}}
\end{center}
\end{document}
